\section{Experiments}
\label{sec:experiments}

\text{TODO: we also need to provide training experiments with ACT and Diffusino policy like in LeVR.}

This section is organised around four research questions, each answered by
one experiment. \textbf{RQ1} (Section~\ref{sec:exp-tracking}): on a 5-DoF
arm, does a position-first IK layer track table-plane strokes better than a
full 6-D pose tracker? \textbf{RQ2} (Section~\ref{sec:exp-wrist}): which
strategy best follows a commanded wrist orientation without wrecking
position? \textbf{RQ3} (Section~\ref{sec:exp-envelope}): how should
infeasible, out-of-envelope targets be handled? \textbf{RQ4}
(Section~\ref{sec:exp-userstudy}): do the objective winners survive contact
with a human operator? RQ1--RQ3 are answered by controlled open-loop
benchmarks; RQ4 is answered by a within-subjects user-study protocol.

All benchmark experiments run in the MuJoCo digital twin (\SI{500}{Hz}
physics, \SI{50}{Hz} control) with scripted mock-hand trajectories that
reproduce the clutch semantics of Eq.~\eqref{eq:clutch}: the trajectory is
defined by a delta curve applied to the EE pose latched at ``grip press.''
Unless noted, metrics aggregate both arms.

\paragraph{Metrics.}
Position error (\(\mathrm{mm}\), target vs.\ measured EE, mean/p95/max) and
its IK-only counterpart (target vs.\ commanded-configuration FK);
orientation error as the geodesic angle between target and measured EE
rotations (\(^\circ\)); \emph{roll lag} (\(\mathrm{ms}\)) as the
cross-correlation peak between commanded and measured roll about the tip
axis; command jerk RMS (\(\mathrm{rad/s^3}\)); peak joint velocity; minimum
joint-limit margin; and solver time. The out-of-envelope suite adds:
time outside, tracking error against the \emph{emitted} (post-policy)
target and against the \emph{raw} operator target while outside, peak joint
velocity while outside, and the re-entry recovery time (raw-target error
\(<\SI{10}{mm}\) sustained \SI{0.2}{s}).

\subsection{RQ1: position-first vs full 6-D tracking}
\label{sec:exp-tracking}

\emph{Does a position-first IK layer track table-plane strokes better than
a full 6-D pose tracker on a 5-DoF arm?} The tracking suite answers this
with three circle radii (3/5/\SI{8}{cm}) and four line-stroke directions on
the table plane, orientation held. Results in
Table~\ref{tab:tracking} (see also the per-trajectory tables in the
supplementary table set). The headline finding of the earlier benchmark holds:
full-cost 6D tracking (\texttt{pink\_full}) pays a large position penalty
on a 5-DoF arm (\SI{34.1}{mm} mean across the full suite vs.\
\SI{15.3}{mm} for \texttt{pink\_relaxed}), while \texttt{scipy\_ls}
(\SI{11.2}{mm}) and the new \texttt{telegrip} split IK (\SI{11.7}{mm})
lead the sweep. \texttt{telegrip}'s position subproblem is 3-DoF and
cannot be disturbed by orientation terms, giving it sub-millimeter
commanded-FK error at a solver cost 13\(\times\) below
\texttt{scipy\_ls} (\SI{0.19}{ms} vs.\ \SI{2.5}{ms} per tick).

\begin{table}[h]
  \centering
  \caption{Tracking suite, representative trajectory
  (\texttt{circle\_r5cm}). Full sweep in the supplementary table set.}
  \label{tab:tracking}
  \small
  \begin{tabular}{lrrrrrrr}
\toprule
method & $\bar{e}_{pos}$ (mm) & $e_{pos}^{95}$ (mm) & $\bar{e}_{ori}$ ($^\circ$) & lag (ms) & jerk (rad/s$^3$) & $\dot{q}_{max}$ & solve (ms) \\
\midrule
pink\_full & 33.7 & 51.7 & 1.99 & -- & 0.33 & 0.48 & 0.22 \\
pink\_relaxed & 17.4 & 23.1 & 9.22 & -- & 0.35 & 0.36 & 0.22 \\
dls & 24.8 & 38.6 & 2.45 & -- & 0.61 & 0.58 & 0.12 \\
mink & 28.3 & 44.3 & 2.01 & -- & 0.68 & 0.60 & 0.19 \\
scipy\_ls & 9.55 & 11.4 & 5.61 & -- & 0.76 & 0.55 & 2.54 \\
telegrip & 9.56 & 11.7 & 7.35 & -- & 0.71 & 0.57 & 0.18 \\
\bottomrule
\end{tabular}

\end{table}

\subsection{RQ2: following a commanded wrist orientation}
\label{sec:exp-wrist}

\emph{Which strategy best follows a commanded wrist orientation without
wrecking position?} The wrist-agility suite holds position at the latch
point while the target orientation oscillates: roll (\(\pm45^\circ\)) and
flex (\(\pm30^\circ\)) at 0.5/1/\SI{2}{Hz}, plus a combined slow-circle +
\SI{1}{Hz}-roll trajectory. This suite quantifies the ``wrist is not
agile'' complaint directly.
Table~\ref{tab:wrist} shows the \SI{1}{Hz} roll oscillation: the armplane
QP tracker lags the roll by \(\approx\)\SI{300}{ms} (about a third of the
oscillation period) with \(28^\circ\) mean orientation error, while the
Telegrip split IK cuts the lag to \(\approx\)\SI{120}{ms} and the error to
\(19^\circ\)~--- at the price of an order of magnitude more command jerk,
as the direct wrist mapping slams into the \SI{3}{rad/s} rate limit that
the \(\pm45^\circ\)@\SI{1}{Hz} profile (peak \SI{4.9}{rad/s}) exceeds by
design.

\begin{table}[h]
  \centering
  \caption{Wrist-agility suite, \SI{1}{Hz} roll oscillation
  (\texttt{wrist\_roll\_f1hz}).}
  \label{tab:wrist}
  \small
  \begin{tabular}{lrrrrrrr}
\toprule
method & $\bar{e}_{pos}$ (mm) & $e_{pos}^{95}$ (mm) & $\bar{e}_{ori}$ ($^\circ$) & lag (ms) & jerk (rad/s$^3$) & $\dot{q}_{max}$ & solve (ms) \\
\midrule
pink\_full & 10.7 & 14.7 & 16.9 & 100.0 & 58.1 & 4.72 & 0.21 \\
pink\_relaxed & 9.84 & 10.8 & 27.8 & 300.0 & 7.40 & 0.61 & 0.22 \\
dls & 9.91 & 11.8 & 22.0 & 140.0 & 211.1 & 3.00 & 0.11 \\
mink & 24.7 & 44.7 & 20.7 & 140.0 & 241.2 & 3.00 & 0.21 \\
scipy\_ls & 9.91 & 10.3 & 11.9 & 80.0 & 61.1 & 4.92 & 2.45 \\
telegrip & 9.91 & 10.7 & 19.5 & 120.0 & 373.3 & 3.00 & 0.19 \\
\bottomrule
\end{tabular}

\end{table}

The thumbstick wrist interface (Section~\ref{sec:thumbstick}) answers the
same complaint differently: it does not chase the commanded orientation at
all, but hands the wrist to the operator. We evaluate it at rehearsal level
only~--- the full teleoperation stack driven by the scripted mock device,
not the controlled benchmark suite~--- because its input is a stick command
the open-loop suite cannot script faithfully. Those rehearsals show its
default hold mode giving the steadiest attitude and the best position
tracking of the wrist patterns, precisely because the attitude is held
rather than tracked, while the decaying soft mode fares worst; the run log
records the sweep. This exposes the limit of the orientation metric here:
for a held attitude the target is ``stay put'' by design, so a low
orientation error against that target is not the same as following the
operator's intent~--- which only the user study (RQ4) can score.

\subsection{RQ3: handling out-of-envelope targets}
\label{sec:exp-envelope}

\emph{How should infeasible targets, driven past the reach boundary, be
handled?} The out-of-envelope suite compares the four policies on three
excursions: \texttt{envelope\_radial} (forward stroke \SI{0.30}{m},
\SI{2}{s} dwell far past \(r_{\max}\), smooth return),
\texttt{envelope\_swoop} (\SI{0.20}{m} dip below the floor), and
\texttt{envelope\_slide} (exit radially, translate laterally while outside,
return)~--- the last distinguishes policies that keep tracking the lateral
component (\texttt{project}, \texttt{slow}) from \texttt{freeze}, which
holds a fixed anchor. Representative results in
Table~\ref{tab:envelope}. With the legacy \texttt{warn} behavior the arm
grinds at its limits \(\approx\)\SI{131}{mm} from its own commanded target;
every active policy keeps the emitted-target error at
\(\approx\)\SI{26}{mm} with bounded joint velocity and \(\le\)\SI{1.8}{s}
recovery. \texttt{project} is the recommended default for data collection:
it minimizes raw-target deviation on the slide trajectory and re-enters
seamlessly.

\begin{table}[h]
  \centering
  \caption{Out-of-envelope suite, \texttt{envelope\_slide} with
  \texttt{pink\_relaxed} (policy comparison; full sweep in the
  supplementary table set).}
  \label{tab:envelope}
  \small
  \begin{tabular}{lrrrrrr}
\toprule
policy & $t_{out}$ (s) & $\bar{e}_{emit}$ (mm) & $\bar{e}_{raw}$ (mm) & $\dot{q}_{max}^{out}$ & $t_{rec}$ (s) & jerk (rad/s$^3$) \\
\midrule
freeze & 6.16 & 24.9 & 156.7 & 0.82 & -- & 3.03 \\
project & 6.16 & 26.9 & 147.7 & 0.87 & -- & 2.99 \\
slow & 6.16 & 26.8 & 148.0 & 0.77 & -- & 2.07 \\
warn & 6.16 & 129.9 & 129.9 & 0.99 & -- & 1.95 \\
\bottomrule
\end{tabular}

\end{table}

\subsection{RQ4: surviving a human operator (user-study protocol)}
\label{sec:exp-userstudy}

\emph{Do the objective winners of RQ1--RQ3 survive contact with a human
operator?} Before that question can be asked of people, the integrated
stack must run as one piece. The full armplane stack (One-Euro filter,
clutch, armplane mapping, envelope policy, IK thread) is exercised against
the digital twin with the scripted mock device: circles (\SI{11.6}{mm} mean
EE tracking with \texttt{pink\_relaxed}, \SI{7.0}{mm} with \texttt{telegrip}),
wrist oscillation, and envelope excursions with each policy. These runs
validate the wiring rather than the methods. The controlled comparison lives
in RQ1--RQ3; the human factors need real operators. The suites above are
open-loop: scripted trajectories cannot
capture how a \emph{human} closes the loop through the headset~--- adapting
to lag,
re-clutching, and reacting to boundary behavior. We therefore pair the
objective benchmark with a formative within-subjects user study
(\(n\approx5\), a pilot scale appropriate for rig-level design decisions).
The full operational runbook, session script, and questionnaire sheets live
in the repository's user-study runbook; this section fixes the
design.

\paragraph{Conditions.}
Four teleoperation configurations, chosen to span the design space the
benchmark identified (Table~\ref{tab:userstudy-conditions}). The
out-of-envelope policy is fixed to \texttt{project} for C1--C3 so the
method is the only varied factor; C4 is the unmodified upstream Telegrip
stack (its own WebXR UI, PyBullet IK, no envelope layer), run via the
native-Telegrip bridge of Section~\ref{sec:system}~--- a \emph{system-level} baseline
rather than an IK swap.

\begin{table}[h]
  \centering
  \caption{User-study conditions.}
  \label{tab:userstudy-conditions}
  \small
  \begin{tabular}{llll}
    \toprule
    & pipeline & IK method & OOE policy \\
    \midrule
    C1 & armplane & \texttt{pink\_full} & \texttt{project} \\
    C2 & armplane & \texttt{pink\_relaxed} & \texttt{project} \\
    C3 & armplane & \texttt{telegrip} (split IK) & \texttt{project} \\
    C4 & upstream Telegrip & PyBullet position IK & none (native) \\
    \bottomrule
  \end{tabular}
\end{table}

\paragraph{Task.}
The primary task is the bimanual \emph{pick--handover--place} used in the
simulation benchmark, transferred to the real rig: the right arm picks a
\SI{40}{mm} soft cube from a marked start pose, hands it to the left arm
over the rig center~--- the region where both arms operate near their
envelope boundaries and coordination is unavoidable~--- and the left arm
places it on a target \SI{15}{cm} away. Three timed trials per condition
(\SI{3}{min} timeout each). A secondary garment task, run once per
condition if session time allows, is a \emph{towel half-fold}: both
grippers grasp adjacent corners of a \(30\times30\)\,cm towel and fold it
onto a marked line (success: \(\ge75\%\) edge overlap by visual check).

\paragraph{Objective measures} (per trial): completion time; success /
failure with failure mode (drop, timeout, unrecoverable pose); number of
re-grips (each grip re-calibration is logged by the IK thread); cumulative
time outside the workspace envelope (envelope margin is in the thread's
periodic debug log; not available in C4); drops; and rig collisions
(observer count). Optionally each trial is recorded as a LeRobot dataset
episode, enabling post-hoc EE path smoothness analysis with the metrics of
this section.

\paragraph{Subjective measures} (after each condition): NASA-TLX raw
workload scales~\citep{nasatlx}; the System Usability
Scale~\citep{sus}; and six custom 7-point Likert items targeting the
qualities this work manipulates~--- perceived positional precision, wrist
responsiveness, predictability at the reach boundary, ease of re-gripping,
freedom of headset placement, and ease of coordinating both arms. After
all conditions: a forced ranking and a semi-structured interview on the
overall feel and process (most/least natural aspects, moments of lost
confidence, strategy changes across conditions).

\paragraph{Procedure.}
Per participant (\(\approx\)75\,min): consent and demographics
(5\,min); guided familiarization on the simulation rehearsal tool
(the simulation rehearsal tool, 10\,min); then the four conditions in a
balanced Latin-square order (the fifth participant repeats the first
order), each comprising \(\sim\)3\,min free practice, three recorded
trials, and the per-condition questionnaires (\(\sim\)12\,min); closing
ranking and interview (10\,min). An experimenter holds the
torque-disable switch throughout; participants may stop at any time.

\paragraph{Analysis.}
With \(n\approx5\), per-condition medians and interquartile ranges are
reported, with Friedman tests (Wilcoxon signed-rank post-hoc) as
exploratory statistics only; interviews are coded thematically. The study
is powered to surface usability regressions and strong preferences, not
small effects~--- its role is to check whether the objective winners of
Sections~\ref{sec:exp-tracking}--\ref{sec:exp-envelope} survive contact
with human operators.
